\uuid{Rw4k}
\titre{Estimation d'un paramètre de la loi de Rayleigh}
\niveau{L3}
\module{Probabilité et statistique}
\chapitre{Estimateur}
\sousChapitre{Intervalle de confiance}
\theme{Loi de Rayleigh, maximum de vraisemblance, biais, théorème central limite, Bienaymé-Tchebychev, loi du chi-deux}
\auteur{}
\datecreate{2026-05-19}
\organisation{}
\difficulte{3}

\contenu{
	
	\texte{
		Soit $\theta > 0$ un paramètre inconnu. On dit qu'une variable aléatoire $X$ suit une \textbf{loi de Rayleigh} de paramètre $\theta$ si elle admet pour densité :
		$$
		f_\theta(x) = \frac{x}{\theta^2} e^{\frac{-x^2}{2\theta^2}} \mathbf{1}_{[0,+\infty[}(x).
		$$
		Soit $X_1, \ldots, X_n$ un échantillon i.i.d. de loi $f_\theta$. On pose $\displaystyle \overline{Q}_n = \frac{1}{n} \sum_{i=1}^n X_i^2.$
	}
	
	\begin{enumerate}
		
		%------- Partie 1 : propriétés de la loi -------
		
%		\item
%		\question{Vérifier que $f_\theta$ est bien une densité de probabilité.}
%		\indication{Effectuer le changement de variable $u = x^2/(2\theta^2)$.}
%		\reponse{
%			On a $f_\theta \geq 0$. Pour $x > 0$, on pose $u = \dfrac{x^2}{2\theta^2}$, soit $du = \dfrac{x}{\theta^2}\,dx$. Alors :
%			\[
%			\int_0^{+\infty} \frac{x}{\theta^2} e^{-x^2/(2\theta^2)}\,dx
%			= \int_0^{+\infty} e^{-u}\,du = 1.
%			\]
%			Donc $f_\theta$ est bien une densité de probabilité.
%		}
		
		\item
		\question{Montrer que $\mathbb{E}[X_1^2] = 2\theta^2$ et en déduire que $\mathrm{Var}(X_1^2) = 4\theta^4$.}
		\indication{Utiliser le même changement de variable $u = x^2/(2\theta^2)$ et la formule $\int_0^{+\infty} u^k e^{-u}\,du = k!$.}
		\reponse{
			Avec $u = \dfrac{x^2}{2\theta^2}$, on a $x^2 = 2\theta^2 u$ et $dx = \dfrac{\theta}{\sqrt{2u}}\,du$, donc :
			\[
			\mathbb{E}[X_1^2]
			= \int_0^{+\infty} x^2 \cdot \frac{x}{\theta^2} e^{-x^2/(2\theta^2)}\,dx
			= \int_0^{+\infty} 2\theta^2 u\, e^{-u}\,du
			= 2\theta^2.
			\]
			De même :
			\[
			\mathbb{E}[X_1^4]
			= \int_0^{+\infty} 4\theta^4 u^2 e^{-u}\,du = 8\theta^4.
			\]
			Donc :
			\[
			\mathrm{Var}(X_1^2) = \mathbb{E}[X_1^4] - \left(\mathbb{E}[X_1^2]\right)^2 = 8\theta^4 - 4\theta^4 = 4\theta^4.
			\]
		}
		
		%------- Partie 2 : estimateur ponctuel -------
		
		\item
		\question{Écrire la log-vraisemblance $\ell(\theta) = \ln L(\theta; x_1, \ldots, x_n)$ de l'échantillon, puis déterminer l'estimateur du maximum de vraisemblance $\hat{\theta}^2_{\mathrm{MV}}$ de $\theta^2$.}
		\indication{Exprimer la log-vraisemblance en fonction de $\theta^2$ puis dériver par rapport à $\theta^2$.}
		\reponse{
			La vraisemblance est :
			\[
			L(\theta; x_1,\ldots,x_n) = \prod_{i=1}^n \frac{x_i}{\theta^2} e^{-x_i^2/(2\theta^2)}.
			\]
			En posant $\varphi = \theta^2$ pour simplifier, la log-vraisemblance vaut :
			\[
			\ell(\varphi) = \sum_{i=1}^n \ln x_i - n\ln\varphi - \frac{1}{2\varphi}\sum_{i=1}^n x_i^2.
			\]
			On dérive par rapport à $\varphi$ :
			\[
			\frac{d\ell}{d\varphi} = -\frac{n}{\varphi} + \frac{1}{2\varphi^2}\sum_{i=1}^n x_i^2.
			\]
			En annulant et en vérifiant que c'est bien un maximum (la dérivée seconde est négative), on obtient :
			\[
			\hat{\theta}^2_{\mathrm{MV}} = \frac{1}{2n}\sum_{i=1}^n X_i^2 = \frac{\overline{Q}_n}{2}.
			\]
		}
		
		\item
		\question{Montrer que $\hat{\theta}^2_n = \hat{\theta}^2_{\mathrm{MV}} = \dfrac{\overline{Q}_n}{2}$ est un estimateur sans biais de $\theta^2$.}
		\reponse{
			Par linéarité de l'espérance et la question 2 :
			\[
			\mathbb{E}[\hat{\theta}^2_n]
			= \frac{1}{2}\,\mathbb{E}[\overline{Q}_n]
			= \frac{1}{2}\,\mathbb{E}[X_1^2]
			= \frac{1}{2} \cdot 2\theta^2 = \theta^2.
			\]
			Donc $\hat{\theta}^2_n$ est sans biais.
		}
		
		\item
		\question{Montrer que $\hat{\theta}^2_n$ converge en moyenne quadratique vers $\theta^2$.}
		\reponse{
			L'estimateur étant sans biais, l'écart quadratique moyen est égal à la variance :
			\[
			\mathrm{EQM}(\hat{\theta}^2_n)
			= \mathrm{Var}(\hat{\theta}^2_n)
			= \frac{1}{4}\,\mathrm{Var}(\overline{Q}_n)
			= \frac{1}{4} \cdot \frac{\mathrm{Var}(X_1^2)}{n}
			= \frac{1}{4} \cdot \frac{4\theta^4}{n}
			= \frac{\theta^4}{n}.
			\]
			Comme $\mathrm{EQM}(\hat{\theta}^2_n) = \dfrac{\theta^4}{n} \xrightarrow[n\to+\infty]{} 0$, l'estimateur $\hat{\theta}^2_n$ converge en moyenne quadratique vers $\theta^2$.
		}
		
		%------- Partie 3 : IC exact -------
		
		\item
		\question{Montrer que pour tout $i$, la variable $Y_i = \dfrac{X_i^2}{\theta^2}$ suit une loi exponentielle de paramètre $\dfrac{1}{2}$.}
		\indication{Calculer la fonction de répartition de $Y_i$ en effectuant le changement de variable $u = x^2/(2\theta^2)$.}
		\reponse{
			Pour $t > 0$ :
			\[
			P(Y_i \leq t)
			= P\!\left(X_i \leq \theta\sqrt{t}\right)
			= \int_0^{\theta\sqrt{t}} \frac{x}{\theta^2} e^{-x^2/(2\theta^2)}\,dx.
			\]
			Avec $u = \dfrac{x^2}{2\theta^2}$, quand $x = \theta\sqrt{t}$ on a $u = \dfrac{t}{2}$, donc :
			\[
			P(Y_i \leq t) = \int_0^{t/2} e^{-u}\,du = 1 - e^{-t/2}.
			\]
			C'est la fonction de répartition d'une loi exponentielle de paramètre $\dfrac{1}{2}$.
		}
		
		\item
		\texte{On admet que la somme de $n$ variables aléatoires indépendantes de loi exponentielle de paramètre $\dfrac{1}{2}$ suit une loi $\chi^2(2n)$.}
		\question{En déduire que $T_n = \dfrac{2n\,\overline{Q}_n}{\theta^2}$ suit une loi $\chi^2(2n)$.}
		\reponse{
			Les variables $Y_i = X_i^2/\theta^2$ sont i.i.d. de loi exponentielle de paramètre $\frac{1}{2}$ d'après la question précédente. Par le résultat admis :
			\[
			T_n = \sum_{i=1}^n Y_i = \frac{1}{\theta^2}\sum_{i=1}^n X_i^2 = \frac{2n\,\overline{Q}_n}{\theta^2} \sim \chi^2(2n).
			\]
		}
		
		\item
		\question{En déduire un intervalle de confiance exact de niveau $95\%$ pour $\theta^2$, symétrique en probabilité. On notera $c_1 = \chi^2_{2n;\,0{,}025}$ et $c_2 = \chi^2_{2n;\,0{,}975}$ les valeurs de la table de la loi $\chi^2(2n)$ aux niveaux $2{,}5\%$ et $97{,}5\%$.}
		\reponse{
			On cherche $c_1$ et $c_2$ tels que $P(c_1 \leq T_n \leq c_2) = 0{,}95$ avec $P(T_n < c_1) = P(T_n > c_2) = 0{,}025$. L'inégalité $c_1 \leq T_n \leq c_2$ est équivalente à :
			\[
			\frac{2n\,\overline{Q}_n}{c_2} \leq \theta^2 \leq \frac{2n\,\overline{Q}_n}{c_1}.
			\]
			L'intervalle de confiance exact de niveau $95\%$ pour $\theta^2$ est donc :
			\[
			I_{0{,}95} = \left[\frac{2n\,\overline{Q}_n}{\chi^2_{2n;\,0{,}975}}\;;\;\frac{2n\,\overline{Q}_n}{\chi^2_{2n;\,0{,}025}}\right].
			\]
		}
		
		%------- Partie 4 : IC asymptotique -------
		
		\item
		\question{À l'aide du théorème central limite, déterminer la loi asymptotique de $\sqrt{n}\!\left(\hat{\theta}^2_n - \theta^2\right)$.}
		\reponse{
			Les variables $X_i^2$ sont i.i.d. d'espérance $2\theta^2$ et de variance $4\theta^4$. Le TCL donne :
			\[
			\sqrt{n}\!\left(\overline{Q}_n - 2\theta^2\right) \xrightarrow[n\to+\infty]{\mathcal{L}} \mathcal{N}(0, 4\theta^4).
			\]
			Comme $\hat{\theta}^2_n = \overline{Q}_n/2$, on obtient :
			\[
			\sqrt{n}\!\left(\hat{\theta}^2_n - \theta^2\right) \xrightarrow[n\to+\infty]{\mathcal{L}} \mathcal{N}(0, \theta^4).
			\]
		}
		
		\item
		\question{En déduire un intervalle de confiance asymptotique de niveau $95\%$ pour $\theta^2$.}
		\reponse{
			D'après la question précédente :
			\[
			\frac{\sqrt{n}\!\left(\hat{\theta}^2_n - \theta^2\right)}{\theta^2} \xrightarrow{\mathcal{L}} \mathcal{N}(0,1).
			\]
			Pour $n$ grand, $P\!\left(\left|\hat{\theta}^2_n - \theta^2\right| \leq \dfrac{1{,}96\,\theta^2}{\sqrt{n}}\right) \approx 0{,}95$.
			En remplaçant $\theta^2$ par $\hat{\theta}^2_n$ dans la marge d'erreur, on obtient l'intervalle asymptotique de niveau $95\%$ :
			\[
			J_{0{,}95} = \left[\hat{\theta}^2_n\!\left(1 - \frac{1{,}96}{\sqrt{n}}\right) \;;\; \hat{\theta}^2_n\!\left(1 + \frac{1{,}96}{\sqrt{n}}\right)\right].
			\]
		}
		
		%------- Partie 5 : Bienaymé-Tchebychev -------
		
		\item
		\question{Sans hypothèse supplémentaire sur la loi de $\hat{\theta}^2_n$, montrer qu'il existe une constante $M > 0$ ne dépendant ni de $\theta$ ni de $n$ telle que, pour tout $\lambda > 0$ :
			$$
			P\!\left(\left|\hat{\theta}^2_n - \theta^2\right| \geq \lambda\theta^2\right) \leq \frac{M}{n}.
			$$
		}
		\indication{Appliquer l'inégalité de Bienaymé-Tchebychev à $\hat{\theta}^2_n$.}
		\reponse{
			L'inégalité de Bienaymé-Tchebychev appliquée à $\hat{\theta}^2_n$ donne, pour tout $\varepsilon > 0$ :
			\[
			P\!\left(\left|\hat{\theta}^2_n - \theta^2\right| \geq \varepsilon\right) \leq \frac{\mathrm{Var}(\hat{\theta}^2_n)}{\varepsilon^2} = \frac{\theta^4}{n\varepsilon^2}.
			\]
			En posant $\varepsilon = \lambda\theta^2$ :
			\[
			P\!\left(\left|\hat{\theta}^2_n - \theta^2\right| \geq \lambda\theta^2\right) \leq \frac{\theta^4}{n\lambda^2\theta^4} = \frac{1}{n\lambda^2}.
			\]
			On pose $M = \dfrac{1}{\lambda^2}$, constante ne dépendant ni de $\theta$ ni de $n$.
		}
		
		\item
		\question{En déduire un intervalle de confiance de niveau au moins $95\%$ pour $\theta^2$ via l'inégalité de Bienaymé-Tchebychev. Comparer sa largeur à celle des intervalles obtenus aux questions 8 et 10 : quel est le prix à payer pour ne pas utiliser la loi exacte de $\hat{\theta}^2_n$ ?}
		\reponse{
			On fixe un niveau de confiance $95\%$, soit $\alpha = 0{,}05$. D'après la question précédente, en prenant $\varepsilon$ tel que $\dfrac{\theta^4}{n\varepsilon^2} \leq 0{,}05$, soit $\varepsilon = \theta^2\sqrt{\dfrac{20}{n}}$, on obtient :
			$$
			P\!\left(\left|\hat{\theta}^2_n - \theta^2\right| \leq \theta^2\sqrt{\frac{20}{n}}\right) \geq 0{,}95.
			$$
			En remplaçant $\theta^2$ par $\hat{\theta}^2_n$, l'intervalle de confiance BT est :
			\[
			K_{0{,}95} = \left[\hat{\theta}^2_n\!\left(1 - \sqrt{\frac{20}{n}}\right) \;;\; \hat{\theta}^2_n\!\left(1 + \sqrt{\frac{20}{n}}\right)\right].
			\]
			Sa demi-largeur relative est $\sqrt{20/n} \approx 4{,}47/\sqrt{n}$, contre $1{,}96/\sqrt{n}$ pour l'IC asymptotique.
			L'IC par Bienaymé-Tchebychev est donc environ $4{,}47/1{,}96 \approx 2{,}3$ fois plus large que l'IC asymptotique : c'est le coût de ne pas exploiter la structure probabiliste précise de l'estimateur.
		}
		
	\end{enumerate}
}